\section{Concept}
Below is a pictures of the poster for our concept idea for a \textbf{Cleaning and Drying station
  for Wet and Dirty Shoes}. Concept Idea  \textbf{Smart Carpet} \autoref{fig:poster}.
%------ First Picture -------------
\begin{figure}[h!]
  \center
  \includegraphics[scale=0.44]{./figures/EN_1.jpg}
  \caption{Low cost low energy concept Idea of a cleaning station, where the user simply A
    rubs their dirty and wet sneakers in the cleaning compartment shown in the figure and dries
    them on the Drying compartment before taking them off and hanging the in the shoe shelf.
    The cleaning compartment has a mixture of  \textbf{water and a cleaning agent underneath}.
    This water is levelled as per the shoe sole height. The user them subs each shoe on the
    brushes soaked into the levelled liquid mixture underneath to remove the dirt and moves to
    the cleaning compartment which contains cleaning re-usable washable cleaning tissues, with a
    cleaning tissue compartment underneath full of clean replacement tissues. The change
    from  dirty to clean Tissues or towels is down automatically.}
  \label{fig:poster}
\end{figure}\\
\subsection{Presentation of Concept Idea}
The procedure used to come up with this concept was through focusing on requirement rather
solely than through solution ideas. This in return created a design that was more realistic and
cost efficient. After having looked at a pretty interesting example of a bad design
(The rimless bicycle) this oriented my focus on requirements more than an ideology. The main
reason why the bicycle idea was a bad design was that, it was solely based on selling the idea
of a futuristic looking bicycle rather than something functional. Since using the requirements
rout would only produce a regular looking bicycle at the end of the day. In making a good
design, we have to take into account the repairability of the product as well as its assembly.
This could also be a huge selling point because it makes the usability aspect of the product
more intuitive.
\subsection{Technical Requirement Specification}
\begin{itemize}
  \item \textbf{Size}: Footprint no larger than  $90cm(\text{length}) \times 30cm(\text{width})
          = 2.7m^2$ height: $20cm$
  \item \textbf{Capacity}: Capable of cleaning \textbf{a pair} of shoe at a time.
  \item \textbf{Material}: Withstand damp shoes, no corrosion, no rot, no rust. The material used should be \textbf{recyclable} and \textbf{durable}.
  \item \textbf{Cleaning}: Brushes and a Solution of \textbf{Water and cleaning Agent}.
  \item \textbf{Humidity}: Keep shoes dry in order to avoid humidity settling in and producing
        unpleasant smells.
  \item \textbf{Cost}: Cost of materials should be kept under 30 Euros.
  \item \textbf{Low Energy Consumption}: Energy consumption should be kept at the minimum.
\end{itemize}

% NOTE: You could take that down in your book or something I like this thought to be honest.
\textbf{There is sunshine for everybody you just have to go outside.}

\subsection{Description}
We are imagining or a product that occupies as much space as a normal shoe rack. Dimensionally
speaking we are talking about a rectangular surface area of roughly $2.7m^2$ and a height of
$20cm/0.02m^2$. with two compartments, one for cleaning and the other for drying shoes and there
is a regular shoe rag for keeping the shoes and letting dry even longer. Refer to the Drawing of the system in the poster \autoref{fig:poster}

\subsection{Functioning of the System}
The system is divided into two compartments namely;  \textbf{Cleaning Compartment} and
\textbf{Drying compartment}
\subsubsection{Cleaning Compartment}
\begin{itemize}
  \item The of area has two compartments, where the use is suppose to stand and physically move
        their fits in a two and fro motion in order to clean each shoe. Inside both
        compartments  we have brushes setup for cleaning the sole of the shoe. Underneath there is the mixture of water and cleaning solution leveled in order to only clean the sole of the shoe(sneakers in this case).
  \item Water and cleaning agent are stored underneath in what we labeled inside the figure as
        cleaning compartment. This is suppose to be manually changed over time as the water
        gets dirty.
\end{itemize}

\subsubsection{Drying Compartment}
\begin{itemize}
  \item The surface of the compartment is similar to that of the cleaning compartment meaning
        there are two shoe compartments for carrying out the same exact task as the cleaning
        compartment with the only difference of what is inside the compartment. We have dry
        reusable towels dry the cleaned shoes.
  \item Underneath  this compartment we also have a bunch of clean dry tissues that
        automatically gets replaced in a pulling motion to uncover the clean tissues for the
        next user.
  \item for shoes that require a quick wipe of the surface, we have included a cleaning
        wrack on the shoe shelf itself.
\end{itemize}

